% Partie : histoire
% Chapitre : biographies
% Section : Kurt Gödel

\documentclass[../../../include/open-logic-section]{subfiles}

\begin{document}

\olfileid[fr]{his}{bio}{god}

\olsection{Kurt G\"odel}

\olphoto{goedel-kurt}{Kurt G{\"o}del}

Kurt G{\"o}del
(\textsc{gueu}-deul)
naquit le
28 avril 1906
à Br{\"u}nn,
dans l'Empire
austro-hongrois
(aujourd'hui
Brno, en
République tchèque).
Curieux et
vif d'esprit,
le jeune Kurtele
était souvent
surnommé dans
sa famille
«~Der kleine
Herr Warum~»
(«~le petit
Monsieur Pourquoi~»).
Il excellait
à l'école
depuis les
classes primaires ;
les mathématiques
furent la seule
matière où il
n'obtint pas
toujours la
meilleure note.
Sa santé fragile
le tenait souvent
absent de l'école,
et il fut dispensé
d'éducation physique.
On lui diagnostiqua
un rhumatisme
articulaire aigu
pendant son enfance.
Toute sa vie,
il crut que
cette maladie
avait laissé
des séquelles
permanentes
sur son cœur,
bien que les
examens médicaux
aient indiqué
le contraire.

G{\"o}del commença
ses études à
l'université
de Vienne
en 1924. Il acheva sa
thèse en~1929 et obtint
son doctorat en~1930.
Il envisageait
d'abord d'étudier
la physique,
mais se tourna
bientôt vers
les mathématiques,
particulièrement
la logique,
en partie sous
l'influence
du philosophe
Rudolf Carnap.
Sa thèse, rédigée
sous la direction
de Hans Hahn,
démontrait le
théorème de
complétude de
la logique
des prédicats
du premier ordre
avec identité
\citep{Godel1929}.
En 1930, il obtint le premier de ses deux résultats
les plus célèbres. Ses premier et second théorèmes
d'incomplétude parurent l'année suivante
\citep{Godel1931}.
Pendant son
séjour à Vienne,
G{\"o}del participa
activement au
Cercle de Vienne,
un groupe de
philosophes
attachés aux
sciences dont
faisait partie
Carnap. Les
résultats de
G{\"o}del influencèrent
particulièrement
les travaux
de ce dernier.

En 1938,
G\"odel épousa
Adele Nimbursky.
Ses parents
n'approuvèrent pas
cette union :
elle avait
six ans de
plus que lui,
était déjà
divorcée et
travaillait
comme danseuse
dans une boîte
de nuit.
Ces pressions
sociales
n'infléchirent
cependant pas
G{\"o}del,
et leur mariage
resta heureux
jusqu'à sa mort.

Après l'annexion
de l'Autriche
par l'Allemagne
nazie en 1938,
G{\"o}del et Adele
émigrèrent en 1940
aux États-Unis,
où il prit
un poste à
l'Institute for
Advanced Study
de Princeton,
dans le
New Jersey.
Malgré son
introversion
et ses
excentricités,
ses années
à Princeton
furent marquées
par des collaborations
fécondes. Il
publia des
essais de
théorie des
ensembles,
de philosophie
et de physique.
Il noua notamment
une amitié
particulièrement
forte avec
son collègue
de l'institut,
Albert Einstein.

Dans ses
dernières années,
la santé
psychique de
G{\"o}del
se dégrada.
L'hospitalisation
de son épouse
en 1977
l'empêcha
de continuer
à préparer
ses repas.
Après des
troubles psychiques
qui l'avaient
accompagné
toute sa vie,
il succomba
à la paranoïa.
Craignant
avec une
intensité
extrême d'être
empoisonné,
G{\"o}del refusa
de manger.
Il mourut
d'inanition
le 14 janvier 1978
à Princeton.

\begin{reading}
Pour une biographie
complète de
G{\"o}del, voir
\citet{Dawson1997}.
Pour d'autres
textes biographiques
et des essais
sur ses contributions
à la logique
et à la philosophie,
voir
\citet{Wang1990},
\citet{Baaz2011},
\citet{Takeuti2003}
et
\citet{Sigmund2007}.

Sa thèse de
doctorat est
disponible
dans l'original
allemand
\citep{Godel1929}.
Le texte
original de
ses théorèmes
d'incomplétude
est publié
dans
\citep{Godel1931}.
Les deux premiers volumes des \emph{Collected Works}
réunissent ses écrits publiés, avec traduction anglaise des textes
allemands \citep{Godel1986,Godel1990}. Le troisième volume
contient une sélection de textes inédits, et les quatrième et
cinquième volumes une sélection de sa correspondance.

Pour une étude
détaillée des
théorèmes
d'incomplétude
de G{\"o}del,
voir
\citet{Smith2013}.
Pour une
présentation
informelle et
philosophique,
voir le podcast
de Mark
Linsenmayer
\citep{Linsenmayer2014}.
\end{reading}

\end{document}
